\documentclass[a4paper,11pt]{article}

% ============================================================
%  Le Collé — Équations différentielles linéaires
%  © 2026 Ahmed Amine Ouqlani
% ============================================================

% ---------- Polices ----------
\usepackage[utf8]{inputenc}
\usepackage[T1]{fontenc}
\usepackage[sc]{mathpazo}          % texte et mathématiques : Palatino
\usepackage[scaled=0.88]{helvet}   % titres et repères : linéale
\linespread{1.07}
\usepackage{amsmath,amssymb,amsthm,mathrsfs,stmaryrd}
\usepackage{microtype}

% ---------- Mise en page ----------
\usepackage[a4paper,top=2.3cm,bottom=2.3cm,left=2.4cm,right=2.4cm,
            headheight=14pt,headsep=0.8cm,footskip=1.1cm]{geometry}
\setlength{\parindent}{0pt}
\setlength{\parskip}{0.45em plus 0.15em minus 0.1em}
\raggedbottom
\emergencystretch=1.5em
\clubpenalty=10000 \widowpenalty=10000 \displaywidowpenalty=10000 \brokenpenalty=10000

% ---------- Couleurs : un bleu de nuit, un ocre, des gris ----------
\usepackage{xcolor}
\definecolor{Ink}{HTML}{1C2533}
\definecolor{Bleu}{HTML}{1F3A5F}
\definecolor{BleuMoyen}{HTML}{4A6A8F}
\definecolor{BleuClair}{HTML}{9DB2C9}
\definecolor{Fond}{HTML}{F2F4F7}
\definecolor{Ocre}{HTML}{B07A26}
\definecolor{Brique}{HTML}{A23B2A}
\definecolor{Gris}{HTML}{5E6672}
\definecolor{Filet}{HTML}{C5CCD6}

% ---------- Packages ----------
\usepackage{enumitem}
\setlist{itemsep=0.25em,topsep=0.35em,parsep=0pt}
\setlist[itemize]{label=\textcolor{BleuMoyen}{\small\textbullet},leftmargin=1.4em}
\setlist[enumerate]{leftmargin=1.8em,label=\textsf{\small\color{Bleu}\arabic*.}}
\usepackage{booktabs,array,multicol}
\usepackage{needspace,float}
\usepackage[font=small,labelfont={sf,color=Bleu},labelsep=period,skip=6pt]{caption}
\usepackage{tikz}
\usetikzlibrary{arrows.meta,calc,decorations.markings}
\usepackage{pgfplots}
\pgfplotsset{compat=1.17}
\usepackage{tkz-tab}
\usepackage[most]{tcolorbox}
\usepackage{titlesec,titletoc}
\usepackage{fancyhdr,lastpage,etoolbox}
\usepackage[colorlinks=true,linkcolor=Bleu,urlcolor=Bleu,citecolor=Bleu,
            pdftitle={Le Collé — Équations différentielles linéaires},pdfauthor={Ahmed Amine Ouqlani}]{hyperref}

\numberwithin{equation}{section}
\numberwithin{figure}{section}
\numberwithin{table}{section}

% ---------- Titres ----------
% numéro du chapitre dans le parcours PT (affiché), distinct de la numérotation interne
\newcommand{\grandnum}{\thesection}\newcommand{\ongletcouleur}{Bleu}
% Chapitres (sections) : grand numéro, titre en linéale, filet ocre
\titleformat{\section}[display]
  {\normalfont\sffamily\color{Ink}}
  {{\small\color{Gris}\textls[200]{\MakeUppercase{Chapitre}}}\\[-0.2em]%
   {\fontsize{64}{64}\selectfont\color{BleuClair}\grandnum}}
  {0.4em}
  {\fontsize{25}{30}\selectfont\bfseries}
  [\vspace{0.5em}{\color{Ocre}\rule{3.2cm}{1.4pt}}]
\titlespacing*{\section}{0pt}{0pt}{1.6em}
\titleformat{name=\section,numberless}[display]
  {\normalfont\sffamily\color{Ink}}{}{0pt}{\fontsize{25}{30}\selectfont\bfseries}
  [\vspace{0.5em}{\color{Ocre}\rule{3.2cm}{1.4pt}}]
\titlespacing*{name=\section,numberless}{0pt}{0pt}{1.6em}

\titleformat{\subsection}
  {\needspace{5\baselineskip}\normalfont\sffamily\Large\bfseries\color{Ink}}
  {\color{Bleu}\thesubsection}{0.8em}{}
\titlespacing*{\subsection}{0pt}{1.8em}{0.5em}
\titleformat{\subsubsection}
  {\needspace{5\baselineskip}\normalfont\sffamily\large\bfseries\color{Ink}}
  {\color{Bleu}\thesubsubsection}{0.7em}{}
\titlespacing*{\subsubsection}{0pt}{1.6em}{0.4em}
\titleformat{\paragraph}[runin]{\normalfont\sffamily\bfseries\color{Ink}}{}{0pt}{}[.]
\titlespacing*{\paragraph}{0pt}{0.9em}{0.6em}
\setcounter{secnumdepth}{3}
\setcounter{tocdepth}{2}

% ---------- Table des matières ----------
\titlecontents{section}[0pt]{\addvspace{0.7em}\sffamily\bfseries\color{Ink}}
  {}{}{\hfill\contentspage}[\addvspace{0.2em}]
\titlecontents{subsection}[2.2em]{\small\linespread{1}\selectfont}
  {\makebox[2.6em][l]{\color{Gris}\thecontentslabel}}{}{\titlerule*[0.7pc]{.}\contentspage}

% ---------- En-têtes et pieds de page ----------
\pagestyle{fancy}
\fancyhf{}
\renewcommand{\sectionmark}[1]{\markboth{#1}{}}
\fancyhead[L]{\footnotesize\sffamily\color{Gris}\leftmark}
\fancyhead[R]{\footnotesize\sffamily\scshape\color{Gris}Le Collé}
\fancyfoot[C]{\small\sffamily\color{Ink}\thepage}
\fancyfoot[L]{\scriptsize\sffamily\hyperlink{toc}{\color{Filet}sommaire}}
\renewcommand{\headrulewidth}{0.4pt}
\renewcommand{\headrule}{{\color{Filet}\hrule height\headrulewidth width\headwidth}}
\renewcommand{\footrulewidth}{0pt}
\fancypagestyle{plain}{\fancyhf{}\fancyfoot[C]{\small\sffamily\color{Ink}\thepage}\renewcommand{\headrulewidth}{0pt}}

% ---------- Énoncés ----------
\newtheoremstyle{enonce}{0pt}{0pt}{}{}{\sffamily\bfseries\color{Bleu}}{.}{0.5em}
  {\thmname{#1}\thmnumber{~#2}\thmnote{\ \normalfont\sffamily\color{Ink}(#3)}}
\newtheoremstyle{exemple}{0.9em}{0.9em}{}{}{\sffamily\bfseries\color{BleuMoyen}}{.}{0.5em}
  {\thmname{#1}\thmnumber{~#2}\thmnote{\ \normalfont\sffamily\color{Ink}(#3)}}
\newtheoremstyle{remarque}{0.7em}{0.7em}{}{}{\sffamily\itshape}{.}{0.5em}{\thmname{#1}\thmnote{\ (#3)}}
\theoremstyle{enonce}
\newtheorem{definition}{Définition}[section]
\newtheorem{theoreme}[definition]{Théorème}
\newtheorem{proposition}[definition]{Proposition}
\newtheorem{corollaire}[definition]{Corollaire}
\theoremstyle{exemple}
\newtheorem{exemple}[definition]{Exemple}
\theoremstyle{remarque}
\newtheorem*{remarque}{Remarque}

\tcbset{
  defi/.style={enhanced,breakable,colback=white,colframe=white,boxrule=0pt,arc=0pt,
     borderline west={1.2pt}{0pt}{Bleu},left=11pt,right=4pt,top=3pt,bottom=3pt,
     before skip=1.1em,after skip=1.1em},
  thm/.style={enhanced,breakable,colback=Fond,colframe=Fond,boxrule=0pt,arc=0pt,
     borderline west={2.4pt}{0pt}{Bleu},left=11pt,right=8pt,top=6pt,bottom=6pt,
     before skip=1.2em,after skip=1.2em}}
\tcolorboxenvironment{definition}{defi}
\tcolorboxenvironment{theoreme}{thm}
\tcolorboxenvironment{proposition}{thm}
\tcolorboxenvironment{corollaire}{thm}

\renewcommand{\proofname}{Démonstration}
\renewcommand{\qedsymbol}{\textcolor{BleuClair}{$\blacksquare$}}
\makeatletter
\renewenvironment{proof}[1][\proofname]{\par\pushQED{\qed}\normalfont
  \topsep0.5em\relax\trivlist\item[\hskip\labelsep\sffamily\itshape\color{Gris}#1.]\ignorespaces}
  {\popQED\endtrivlist\@endpefalse}
\makeatother

% ---------- Méthode : filets haut et bas, titre en bandeau ----------
\makeatletter\newcommand{\methlabelfix}{\protected@edef\@currentlabel{\thetcbcounter}}\makeatother
\newtcolorbox[auto counter,number within=section]{methode}[2][]{%
  enhanced,unbreakable,colback=white,colframe=white,boxrule=0pt,arc=0pt,
  borderline north={0.9pt}{0pt}{Bleu},borderline south={0.9pt}{0pt}{Bleu},
  left=4pt,right=4pt,top=8pt,bottom=8pt,before skip=1.4em,after skip=1.4em,
  before upper={\methlabelfix\phantomsection\addcontentsline{met}{metentry}{\protect\makebox[2.6em][l]{\thetcbcounter}#2}%
     {\sffamily\bfseries\small\color{white}\colorbox{Bleu}{\strut\ MÉTHODE~\thetcbcounter\ }}%
     \ \ {\sffamily\bfseries\color{Bleu}#2}\par\smallskip},#1}

% ---------- Repères pédagogiques : un filet de couleur, un mot-clé ----------
\newcommand{\definirbloc}[3]{% nom, couleur, libellé
  \newtcolorbox{#1}[1][]{enhanced,breakable,blanker,
    borderline west={1.4pt}{0pt}{#2},left=10pt,right=2pt,top=1pt,bottom=1pt,
    before skip=1em,after skip=1em,
    before upper={{\sffamily\bfseries\footnotesize\color{#2}\textls[120]{\MakeUppercase{#3}}}\quad},##1}}
\definirbloc{reflexe}{Bleu}{Réflexe}
\definirbloc{astuce}{Ocre}{Astuce}
\definirbloc{attention}{Brique}{Attention}
\definirbloc{enpratique}{BleuMoyen}{En pratique}
\definirbloc{plusloin}{Gris}{Pour aller plus loin}
\definirbloc{strategie}{Bleu}{Stratégie}
\definirbloc{lecture}{BleuMoyen}{Lire le schéma}
% compatibilité avec les parties reprises
\newenvironment{garde}{\begin{attention}}{\end{attention}}
\newenvironment{complementHP}{\begin{plusloin}}{\end{plusloin}}
\newenvironment{interpretation}{\begin{enpratique}}{\end{enpratique}}

% ---------- Exercices, corrigés ----------
\newcounter{exo}[section]
\renewcommand{\theexo}{\thesection.\arabic{exo}}
\newcommand{\niveau}[1]{{\color{Ocre}\footnotesize
  \ifcase#1\or$\bullet\circ\circ\circ\circ$\or$\bullet\bullet\circ\circ\circ$\or$\bullet\bullet\bullet\circ\circ$%
  \or$\bullet\bullet\bullet\bullet\circ$\or$\bullet\bullet\bullet\bullet\bullet$\fi}}
\newcommand{\serie}[1]{\par\needspace{8\baselineskip}\vspace{1.5em}%
  \noindent{\sffamily\bfseries\color{Bleu}#1}\par\nobreak\vspace{-0.4em}%
  \noindent{\color{Filet}\rule{\linewidth}{0.4pt}}\par\nobreak\vspace{0.1em}}
\newenvironment{exo}[3]{%
  \par\needspace{5\baselineskip}\vspace{0.75em}\refstepcounter{exo}\label{e:#1}%
  \def\exolab{#1}%
  \noindent{\sffamily\bfseries\color{Bleu}Exercice~\theexo}\quad\niveau{#2}%
  \hfill{\footnotesize\sffamily\color{Gris}#3}\par\nobreak\vspace{0.15em}}%
  {\par\nopagebreak\noindent\hfill{\scriptsize\sffamily\hyperref[c:\exolab]{\color{Gris}corrigé p.~\pageref{c:\exolab}}}\par}
\newenvironment{corrige}[1]{%
  \par\needspace{6\baselineskip}\vspace{1em}\phantomsection\label{c:#1}%
  \noindent{\sffamily\bfseries\color{Bleu}Exercice~\ref{e:#1}}\quad
  {\scriptsize\sffamily\hyperref[e:#1]{\color{Gris}énoncé p.~\pageref{e:#1}}}\par\nobreak\vspace{0.2em}}%
  {\hfill\textcolor{BleuClair}{$\square$}\par}
\newcommand{\etape}[1]{\par\smallskip\noindent{\sffamily\bfseries\small\color{BleuMoyen}#1.}\ }
\newenvironment{solution}{\par\noindent{\sffamily\bfseries\small\color{BleuMoyen}Solution.}\ }{\hfill\textcolor{BleuClair}{$\square$}\par}
\newcommand{\chercher}{\par\smallskip\noindent{\color{Filet}\rule[0.5ex]{1.2cm}{0.4pt}}\
  {\footnotesize\sffamily\itshape\color{Gris}laisser chercher l'étudiant avant de poursuivre}\par\smallskip}

% ---------- Guide de séance ----------

\newcommand{\marge}[1]{\leavevmode\marginpar{\raggedleft\scriptsize\sffamily\color{Gris}\textls[100]{\MakeUppercase{#1}}}\ignorespaces}
\newcommand{\Etape}[4]{%
  \subsection[Étape #1 --- #3 {\color{Gris}(#2)}]{#3}\label{etape:#1}%
  \noindent{\small\sffamily\color{Gris}{\color{Ocre}\bfseries Étape #1}\quad·\quad #2\quad·\quad #4}\par\medskip}
\newcommand{\question}[1]{\par\smallskip\noindent{\sffamily\bfseries\small\color{BleuMoyen}Question.}\ \textit{#1}\par\smallskip}
\newcommand{\rappel}[1]{\par\smallskip\noindent\phantomsection\addcontentsline{rap}{rapentry}{#1}%
  {\sffamily\bfseries\scriptsize\colorbox{Fond}{\color{Gris}\strut RAPPEL}}\ {\sffamily\small\color{Gris}#1.}\ }
\newcommand{\pourquoi}{\par\noindent{\sffamily\bfseries\small\color{Ocre}Pourquoi ?}\ }

% Ouverture de chapitre : \ouverture{objectifs}{prérequis}{durée}
\newcommand{\ouverture}[3]{%
  \begin{tcolorbox}[enhanced,colback=Fond,colframe=Fond,arc=0pt,boxrule=0pt,left=12pt,right=12pt,top=10pt,bottom=10pt,
     before skip=1.2em,after skip=1.6em]
  \begin{minipage}[t]{0.56\linewidth}
  {\sffamily\bfseries\small\color{Bleu}OBJECTIFS}\par\vspace{0.3em}\small #1
  \end{minipage}\hfill
  \begin{minipage}[t]{0.38\linewidth}
  {\sffamily\bfseries\small\color{Bleu}PRÉREQUIS}\par\vspace{0.3em}\small #2
  \ifx\relax#3\relax\else\par\vspace{0.6em}{\sffamily\bfseries\small\color{Bleu}DURÉE}\par\vspace{0.2em}\small #3\fi
  \end{minipage}
  \end{tcolorbox}}
% Fin de chapitre
\newenvironment{essentiel}{\par\needspace{12\baselineskip}\vspace{1.5em}%
  \begin{tcolorbox}[enhanced,unbreakable,colback=white,colframe=white,boxrule=0pt,arc=0pt,
   borderline north={1.4pt}{0pt}{Ocre},left=2pt,right=2pt,top=8pt,bottom=4pt]
  {\sffamily\bfseries\color{Ocre}L'ESSENTIEL DU CHAPITRE}\par\medskip}{\end{tcolorbox}}

% ---------- Notations ----------
\newcommand{\K}{\mathbb{K}}
\newcommand{\R}{\mathbb{R}}
\newcommand{\C}{\mathbb{C}}
\newcommand{\N}{\mathbb{N}}
\newcommand{\Z}{\mathbb{Z}}
\newcommand{\Mn}[1]{\mathcal{M}_{n}(#1)}
\newcommand{\GLn}[1]{\mathrm{GL}_n(#1)}
\newcommand{\LE}{\mathcal{L}(E)}
\newcommand{\rg}{\operatorname{rg}}
\newcommand{\vect}{\operatorname{Vect}}
\newcommand{\Mat}{\operatorname{Mat}}
\newcommand{\Id}{\operatorname{Id}}
\newcommand{\sh}{\operatorname{sh}}
\newcommand{\ch}{\operatorname{ch}}
\renewcommand{\th}{\operatorname{th}}
\DeclareMathOperator{\argsh}{argsh}
\newcommand{\vv}[1]{\overrightarrow{#1}}
\newcommand{\ve}[1]{\vec{#1}}
\newcommand{\ient}[2]{\llbracket #1,#2\rrbracket}
\newcommand{\dt}[1]{\left|\begin{matrix}#1\end{matrix}\right|}
\newcommand{\dis}{\displaystyle}

% ---------- Figures ----------
\pgfplotsset{trig format plots=rad}
\pgfplotsset{courbe/.style={axis lines=middle,axis equal image,
  axis line style={Gris,-{Stealth[length=5pt]}},
  xlabel={$x$},ylabel={$y$},ticklabel style={font=\scriptsize\sffamily,color=Gris},
  every axis x label/.style={at={(ticklabel* cs:1)},anchor=west,font=\small},
  every axis y label/.style={at={(ticklabel* cs:1)},anchor=south,font=\small},
  samples=300,smooth,width=0.6\linewidth}}
\tikzset{arcetude/.style={Bleu,line width=1.25pt},
         reste/.style={BleuClair,line width=1pt},
         tgte/.style={Ocre,line width=0.9pt},
         asy/.style={Gris,dashed,line width=0.6pt},
         pt/.style={circle,fill=Ink,inner sep=1.2pt}}

% ---------- Navigation ----------
\makeatletter
\newcommand*\l@rapentry{\@dottedtocline{1}{0em}{0em}}
\newcommand*\l@metentry{\@dottedtocline{1}{0em}{0em}}
\newcommand{\listemethodes}{\@starttoc{met}}
\newcommand{\listerappels}{\@starttoc{rap}}
\makeatother
% onglets de chapitre sur le bord droit
\usepackage{eso-pic}
\newif\ifonglets \ongletsfalse
% ---------------------------------------------------------------------------------------------
% NIVEAU DU CHAPITRE. \PTfalse : programme de première année (PTSI) seulement, version actuelle du site.
% \PTtrue : version complète de PT (second ordre à coefficients continus, recollement, variation de la
% constante au second ordre, séries entières). Rien n'est supprimé : le contenu de PT est dans les blocs \ifPT ... \fi.
% Pour le réintroduire : remplacer \PTfalse par \PTtrue (même ligne dans la présentation, la fiche,
% l'exercice bilan) et passer PT = True dans outils/td/meta.py et outils/td/selection.py.
\newif\ifPT \ifdefined\VERSIONPT \PTtrue \else \PTfalse \fi
% ---------------------------------------------------------------------------------------------
\newcommand{\dessineronglet}{%
  \edef\ongcol{\ongletcouleur}%
  \edef\ongnum{\grandnum}%
  \pgfmathsetmacro{\ongy}{3.2+2.3*(\ongnum-1)}%
  \begin{tikzpicture}[remember picture,overlay]
  \fill[\ongcol] ([xshift=-0.75cm,yshift=-\ongy cm]current page.north east)
     rectangle ([yshift=-\ongy cm-1.9cm]current page.north east);
  \node[text=white,font=\sffamily\bfseries\large] at ([xshift=-0.37cm,yshift=-\ongy cm-0.95cm]current page.north east)
     {\ongnum};
  \end{tikzpicture}}
\AddToShipoutPictureBG{\ifonglets\ifnum\value{section}>0 \dessineronglet\fi\fi}
% blocs de navigation et d'enseignement
\definirbloc{commencer}{Gris}{Si vous commencez ici}
\definirbloc{horsprogramme}{Brique}{Hors programme}
\newcommand{\HP}{\,{\sffamily\bfseries\scriptsize\textcolor{Brique}{[HP]}}}
\newcommand{\source}[1]{{\footnotesize\sffamily\color{Gris}\,(#1)}}
\newcommand{\coupdepouce}[1]{\par\vspace{-0.2em}\noindent{\footnotesize\sffamily\color{Gris}%
  \textbf{Coup de pouce} (à donner si l'étudiant bloque) : #1}\par}
\newcommand{\pointarret}[1]{\par\vspace{1em}\noindent{\color{Filet}\rule[0.5ex]{\linewidth}{0.4pt}}\par\nobreak
  \noindent{\footnotesize\sffamily\color{Ocre}\bfseries POINT D'ARRÊT POSSIBLE}\quad{\footnotesize\sffamily\color{Gris}#1}\par\vspace{0.6em}}
\newcommand{\demarche}[1]{\noindent{\sffamily\bfseries\small\color{BleuMoyen}Démarche.}\ \textit{#1}\par\smallskip}

% ---------- Corrigés : idées, astuces, pièges, à retenir ----------
\newcommand{\choix}[1]{\par\noindent{\sffamily\bfseries\small\color{BleuMoyen}Choix de la méthode.}\ #1\par\smallskip}
\newcommand{\gain}[1]{\par\noindent{\sffamily\bfseries\small\color{Ocre}Astuce.}\ #1\par}
\newcommand{\piege}[1]{\par\noindent{\sffamily\bfseries\small\color{Brique}Piège.}\ #1\par}
\newtcolorbox{retenirbox}{enhanced,breakable,colback=Fond,colframe=Fond,arc=0pt,boxrule=0pt,
  borderline west={2pt}{0pt}{Ocre},left=8pt,right=6pt,top=4pt,bottom=4pt,before skip=0.7em,after skip=0.5em}
\newcommand{\retenir}[1]{\begin{retenirbox}{\sffamily\bfseries\footnotesize\color{Ocre}À RETENIR}\quad #1\end{retenirbox}}
\newcommand{\TD}[1]{{\color{Ocre}\bfseries TD\,#1}\ ·\ }
\newcommand{\q}[1]{\par\smallskip\noindent{\sffamily\bfseries\color{Bleu}#1}\ }
\providecommand{\iddots}{\mathinner{\mkern1mu\raise1pt\hbox{.}\mkern2mu\raise4pt\hbox{.}\mkern2mu\raise7pt\hbox{.}\mkern1mu}}


\begin{document}
\renewcommand{\grandnum}{4}
\pagenumbering{roman}
\begin{titlepage}
\thispagestyle{empty}
\begin{tikzpicture}[remember picture,overlay]
  \fill[Bleu] (current page.north west) rectangle ([yshift=-15.5cm]current page.north east);
  \node[anchor=north west,text=white,font=\sffamily\small] at ([xshift=2.6cm,yshift=-3.2cm]current page.north west)
     {\textls[250]{MATHÉMATIQUES\ \ ·\ \ CPGE PT\ \ ·\ \ DEUXIÈME ANNÉE}};
  \node[anchor=north west,text=white] at ([xshift=2.5cm,yshift=-4.4cm]current page.north west)
     {\fontsize{62}{66}\selectfont\scshape Le Collé};
  \draw[Ocre,line width=1.6pt] ([xshift=2.6cm,yshift=-7.3cm]current page.north west) -- ++(4.2cm,0);
  \node[anchor=north west,text=white,font=\sffamily\Large\bfseries,align=left]
     at ([xshift=2.6cm,yshift=-8.0cm]current page.north west) {Équations différentielles};
  \node[anchor=north west,text=white!90!Bleu,font=\sffamily\large,align=left]
     at ([xshift=2.6cm,yshift=-9.0cm]current page.north west) {Équations linéaires du premier et du second ordre\\[0.15cm]\ifPT Coefficients constants et coefficients continus\else Coefficients constants · programme de première année (PTSI)\fi};
  \node[anchor=north west,text=Ink,font=\sffamily\small,align=left,text width=12cm]
     at ([xshift=2.6cm,yshift=-17.4cm]current page.north west)
     {\textcolor{Bleu}{\bfseries Cours, méthodes et exercices corrigés}\\[0.3cm]\textcolor{Gris}{Le premier chapitre de l'année : la structure des solutions, les méthodes de résolution, et le théorème de Cauchy linéaire.}};
  \node[anchor=south west,text=Ink,font=\sffamily\large] at ([xshift=2.6cm,yshift=3.0cm]current page.south west) {Ahmed Amine Ouqlani};
  \draw[Filet] ([xshift=2.6cm,yshift=2.6cm]current page.south west) -- ([xshift=-2.4cm,yshift=2.6cm]current page.south east);
\end{tikzpicture}
\end{titlepage}
\thispagestyle{plain}
\section*{Avant-propos}
\markboth{Avant-propos}{}
Ce livre accompagne la séance sur les équations différentielles linéaires. Il reprend, dans l'ordre et avec les mêmes
notations, la présentation projetée en séance, et sert de référence pour travailler seul ensuite : chaque résultat y est
énoncé avec ses hypothèses, chaque méthode est illustrée par un exemple rédigé, et chaque partie se termine par quelques
exercices, corrigés à la fin du livret.

\ifPT
Le programme de PT prolonge la première année : les équations du premier ordre et du second ordre à coefficients constants
sont supposées connues, mais elles sont reprises ici en entier, car elles servent dans toutes les autres méthodes.
La nouveauté de deuxième année est l'équation du second ordre à \emph{coefficients continus}
$y''+a(t)y'+b(t)y=c(t)$ : théorème de Cauchy linéaire, structure de l'ensemble des solutions, et quelques méthodes de
résolution guidées (solution polynomiale, solution connue ne s'annulant pas, changement de variable, séries entières).

La recherche de solutions développables en série entière (§1.4) suppose connu le chapitre sur les séries entières ;
si ce n'est pas encore le cas, on peut la laisser de côté sans gêner la suite. Les notions hors programme sont signalées \HP.
Les compléments (TD complet et ses corrections, sélection de cinq problèmes, fiche de référence) suivent la même numérotation.

\else
Cette version du livret suit le \emph{programme de première année} (PTSI) : équations du premier ordre $y'+a(t)y=b(t)$ à
coefficients continus, et équations du second ordre à \emph{coefficients constants} $y''+ay'+by=c(t)$. C'est la base sur
laquelle s'appuie le cours de deuxième année ; elle est reprise ici en entier, avec la structure de l'ensemble des solutions
(espaces vectoriels, dimension), le théorème de Cauchy et les méthodes de recherche d'une solution particulière.

Les compléments de PT (second ordre à coefficients continus, recollement, variation de la constante au second ordre,
solutions développables en série entière) seront ajoutés quand ils auront été vus en classe. Les notions hors programme sont
signalées \HP. Les compléments (TD complet et ses corrections, sélection de cinq problèmes, fiche de référence) suivent la
même numérotation.

\fi
\vfill
{\footnotesize\sffamily\color{Gris}\textcopyright~2026 Ahmed Amine Ouqlani. Tous droits réservés.\par}
\clearpage
\thispagestyle{plain}
\hypertarget{toc}{}
\renewcommand{\contentsname}{Sommaire}
\tableofcontents
\clearpage
\pagenumbering{arabic}
\ongletstrue
% =====================================================================
%            CHAPITRE 1 — ÉQUATIONS DIFFÉRENTIELLES LINÉAIRES
% =====================================================================
\clearpage
\section{Équations différentielles linéaires}\label{sec:ed}

Une équation différentielle relie une fonction inconnue $y$ à ses dérivées. Lorsqu'elle est \emph{linéaire}, l'ensemble de
ses solutions a une structure très simple : une solution particulière, plus toutes les solutions de l'équation sans second
membre, qui forment un espace vectoriel de dimension $1$ (premier ordre) ou $2$ (second ordre). Tout le chapitre consiste à
exploiter cette structure : trouver l'espace des solutions homogènes, puis une seule solution particulière.

\ifPT
\ouverture{%
\begin{itemize}[leftmargin=1.1em,itemsep=1pt]
\item Connaître la structure de l'ensemble des solutions et le théorème de Cauchy linéaire.
\item Résoudre les équations du premier ordre et du second ordre à coefficients constants.
\item Résoudre une équation du second ordre à coefficients continus en suivant une méthode indiquée.
\item Recoller des solutions en un point où le coefficient dominant s'annule.
\end{itemize}}{%
Dérivation, primitives usuelles, intégration par parties ; nombres complexes, exponentielle complexe ; espaces vectoriels
(sous-espace, dimension) ; pour le §1.4, séries entières.}{Deux séances}
\else
\ouverture{%
\begin{itemize}[leftmargin=1.1em,itemsep=1pt]
\item Connaître la structure de l'ensemble des solutions et le théorème de Cauchy.
\item Résoudre une équation du premier ordre : équation homogène, variation de la constante.
\item Résoudre une équation du second ordre à coefficients constants pour les seconds membres usuels.
\item Utiliser le principe de superposition et résoudre un problème de Cauchy.
\end{itemize}}{%
Dérivation, primitives usuelles, intégration par parties ; nombres complexes, exponentielle complexe ; espaces vectoriels
(sous-espace, dimension).}{Une séance}
\fi

Dans tout le chapitre, $I$ est un intervalle de $\R$ non réduit à un point et $\K=\R$ ou $\C$.

% ---------------------------------------------------------------------
\subsection{Vocabulaire et structure des solutions}\label{ss:structure}

\begin{definition}
\ifPT
Une \emph{équation différentielle linéaire scalaire} d'ordre $1$, resp.\ $2$, est une équation de la forme
\[(E)\quad y'+a(t)y=b(t),\qquad\text{resp.}\qquad (E)\quad y''+a(t)y'+b(t)y=c(t),\]
où $a$, $b$ (et $c$) sont des fonctions \emph{continues} de $I$ dans $\K$.
\else
Une \emph{équation différentielle linéaire} du premier ordre, resp.\ du second ordre à coefficients constants, est une
équation de la forme
\[(E)\quad y'+a(t)y=b(t),\qquad\text{resp.}\qquad (E)\quad y''+ay'+by=c(t),\]
où $a$, $b$ sont des fonctions \emph{continues} de $I$ dans $\K$ au premier ordre, et où, au second ordre, $a,b\in\K$ sont
des constantes et $c$ est une fonction continue de $I$ dans $\K$.
\fi Une \emph{solution} est une fonction $y$ dérivable
(resp.\ deux fois dérivable) sur $I$ qui vérifie l'égalité en tout point de $I$. Le \emph{second membre} est $b$ (resp.\ $c$) ;
l'\emph{équation homogène associée} $(H)$ est obtenue en le remplaçant par $0$.
\end{definition}

On rencontre souvent l'équation sous la forme $\alpha(t)y'+\beta(t)y=\gamma(t)$\ifPT, ou $\alpha(t)y''+\dots$\fi{} : on ne peut la mettre
sous la forme ci-dessus (forme \emph{normalisée}) qu'en divisant par $\alpha(t)$, donc \emph{sur un intervalle où $\alpha$ ne
s'annule pas}. \ifPT Les points où $\alpha$ s'annule sont traités à part (§\ref{ss:recollement}).\else On résout alors
l'équation sur chacun des intervalles où $\alpha$ ne s'annule pas.\fi

\begin{proposition}[structure]\label{p:structure}
Soit \ifPT$L(y)=y''+a(t)y'+b(t)y$\else$L(y)=y''+ay'+by$\fi{} (ou $L(y)=y'+a(t)y$ au premier ordre).
\begin{enumerate}
\item $L$ est linéaire : $L(\lambda y+\mu z)=\lambda L(y)+\mu L(z)$. L'ensemble $S_H$ des solutions de $(H)$ est donc un
\emph{espace vectoriel} (le noyau de $L$).
\item \emph{Principe de superposition.} Si $L(y_1)=c_1$ et $L(y_2)=c_2$, alors $L(y_1+y_2)=c_1+c_2$.
\item Si $y_p$ est une solution de $(E)$, les solutions de $(E)$ sont exactement les fonctions $y_p+h$, avec $h\in S_H$ :
\[S_E=y_p+S_H.\]
\end{enumerate}
\end{proposition}

\begin{proof}
Le 1 vient de la linéarité de la dérivation ; le 2 en est un cas particulier. Pour le 3 : $L(y)=c=L(y_p)$ équivaut à
$L(y-y_p)=0$, c'est-à-dire $y-y_p\in S_H$.
\end{proof}

\begin{reflexe}
Résoudre une équation linéaire, c'est toujours faire deux choses séparées : \textbf{(1)} résoudre l'équation homogène ;
\textbf{(2)} trouver \emph{une} solution particulière. Le principe de superposition permet de découper un second membre compliqué
en morceaux simples.
\end{reflexe}

\begin{theoreme}[Cauchy linéaire]\label{t:cauchy}
Soit $t_0\in I$.
\begin{itemize}
\item Ordre $1$ : pour tout $y_0\in\K$, il existe une unique solution de $y'+a(t)y=b(t)$ sur $I$ telle que $y(t_0)=y_0$.
\item Ordre $2$ : pour tous $y_0,y_1\in\K$, il existe une unique solution de \ifPT$y''+a(t)y'+b(t)y=c(t)$\else$y''+ay'+by=c(t)$\fi{} sur $I$ telle que
$y(t_0)=y_0$ et $y'(t_0)=y_1$.
\end{itemize}
\end{theoreme}

La démonstration au second ordre est hors programme ; au premier ordre, elle découle du théorème~\ref{t:ordre1}. Deux
conséquences reviennent constamment.

\begin{corollaire}\label{c:dim}
\begin{enumerate}
\item Une solution de l'équation homogène qui s'annule en un point (au premier ordre), ou qui s'annule avec sa dérivée en un
point (au second ordre), est la fonction nulle.
\item $S_H$ est une droite vectorielle au premier ordre, un plan vectoriel au second ordre : l'application
$h\mapsto\big(h(t_0),h'(t_0)\big)$ est un isomorphisme de $S_H$ sur $\K^2$.
\end{enumerate}
\end{corollaire}

\begin{proof}
1. La fonction nulle est solution du même problème de Cauchy ; par unicité, c'est la seule.
2. Au second ordre, $\Phi:h\mapsto(h(t_0),h'(t_0))$ est linéaire ; elle est surjective (existence) et injective (point 1). Donc
$\dim S_H=\dim\K^2=2$. Même raisonnement au premier ordre avec $h\mapsto h(t_0)$.
\end{proof}

\begin{enpratique}
Pour montrer que deux solutions $y_1,y_2$ de $(H)$ engendrent $S_H$ au second ordre, il suffit de vérifier qu'elles ne sont pas
proportionnelles : elles forment alors une famille libre de deux vecteurs dans un plan. Si on ne voit pas qu'elles ne sont pas
proportionnelles, on calcule en un point $t_0$ le déterminant $\dt{y_1(t_0)&y_2(t_0)\\y_1'(t_0)&y_2'(t_0)}$ : il est non nul
si et seulement si les vecteurs $\Phi(y_1)$ et $\Phi(y_2)$ de $\K^2$ forment une base (chapitre \emph{Déterminants}).
\end{enpratique}


% ---------------------------------------------------------------------
\subsection{Équations du premier ordre}\label{ss:ordre1}

\begin{theoreme}[équation homogène]\label{t:ordre1}
Soit $a$ continue sur $I$ et $A$ une primitive de $a$ sur $I$. Les solutions de $y'+a(t)y=0$ sur $I$ sont les fonctions
\[t\mapsto Ke^{-A(t)},\qquad K\in\K.\]
\end{theoreme}

\begin{proof}
Si $y$ est solution, $\big(y\,e^{A}\big)'=(y'+ay)e^{A}=0$ sur l'intervalle $I$, donc $y\,e^{A}$ est constante. Réciproquement,
$Ke^{-A}$ est solution.
\end{proof}

\begin{attention}
On écrit l'équation sous forme \emph{normalisée} $y'+a(t)y=b(t)$ avant de prendre une primitive : pour $ty'+y=\ln t$ sur
$\R_+^*$, c'est $a(t)=\frac1t$, et non $a(t)=1$. Le signe compte : $y'=a(t)y$ donne $Ke^{+A(t)}$.
\end{attention}

\begin{methode}{Trouver une solution particulière au premier ordre}\label{m:part1}
\begin{enumerate}
\item \emph{Solution évidente :} une constante, un polynôme de bas degré\dots\ (essayer $y=$ constante si le second membre est constant).
\item \emph{Coefficient $a$ constant et second membre $P(t)e^{mt}$} ($P$ polynôme, $m\in\K$) : chercher $y=Q(t)e^{mt}$ avec $Q$
polynôme de même degré que $P$ si $m\neq-a$, de degré $\deg P+1$ si $m=-a$ (résonance).
Pour $\cos(\omega t)$ ou $\sin(\omega t)$ : résoudre avec $e^{i\omega t}$, puis prendre la partie réelle ou imaginaire
(les coefficients doivent être réels).
\item \emph{Variation de la constante :} chercher $y=\lambda(t)e^{-A(t)}$. En reportant, les termes en $\lambda$ disparaissent et il reste
\[\lambda'(t)=b(t)\,e^{A(t)},\]
qu'on intègre. Méthode toujours applicable, à garder quand rien de simple n'apparaît.
\end{enumerate}
\end{methode}

\begin{exemple}[second membre polynomial]
$y'-y=t^2-3t+1$. Ici $a=-1$ et $m=0\neq1$ : on cherche $y=\alpha t^2+\beta t+\gamma$. En identifiant
$(2\alpha t+\beta)-(\alpha t^2+\beta t+\gamma)=t^2-3t+1$ : $\alpha=-1$, $\beta=1$, $\gamma=0$. Solutions : $y=-t^2+t+Ke^t$.
\end{exemple}

\begin{exemple}[second membre trigonométrique]
$y'-3y=-7\sin(4t)$. On résout $z'-3z=-7e^{4it}$ en cherchant $z=\nu e^{4it}$ : $(4i-3)\nu=-7$, $\nu=\frac{-7}{-3+4i}=\frac{21+28i}{25}$.
La partie imaginaire de $z$ est solution de l'équation initiale (dont les coefficients sont réels) :
$y_p=\operatorname{Im}\!\big(\tfrac{21+28i}{25}(\cos4t+i\sin4t)\big)=\tfrac{28}{25}\cos4t+\tfrac{21}{25}\sin4t$. Solutions :
$y=y_p+Ke^{3t}$.
\end{exemple}

\begin{exemple}[variation de la constante]
$ty'+y=\ln t$ sur $\R_+^*$, soit $y'+\frac1ty=\frac{\ln t}t$. Homogène : $y=\frac Kt$. On pose $y=\frac{\lambda(t)}t$ : $\lambda'=\ln t$,
$\lambda=t\ln t-t$. Solutions : $y=\ln t-1+\frac Kt$.
\end{exemple}

\subsubsection*{Équations intégrales}

Une équation où l'inconnue figure sous une intégrale se ramène, en dérivant, à un problème de Cauchy du premier ordre : c'est
une application directe de ce qui précède.

\begin{exemple}[une équation intégrale]
Cherchons les fonctions continues $f$ telles que $f(x)=1+\int_0^xtf(t)\,dt$ pour tout réel $x$. Le membre de droite est
dérivable (intégrale de fonction continue) ; donc $f$ l'est, et $f'(x)=xf(x)$, avec $f(0)=1$. Réciproquement, une solution de ce
problème de Cauchy vérifie l'équation intégrale (intégrer $f'=xf$ de $0$ à $x$). Par le théorème~\ref{t:ordre1}, $f(x)=e^{x^2/2}$.
\end{exemple}

\begin{methode}{Équation intégrale}\label{m:integrale}
Si l'inconnue apparaît sous une intégrale $\int_{x_0}^x$, dériver pour obtenir une équation différentielle, et relever la
condition initiale en $x=x_0$. Conclure par équivalence (ou vérifier que la solution trouvée convient).
\end{methode}

\ifPT
\subsubsection*{Recollement}\label{ss:recollement}

\begin{methode}{Recoller des solutions}\label{m:recollement}
Pour $\alpha(t)y'+\beta(t)y=\gamma(t)$ quand $\alpha$ s'annule en $t_0$ :
\begin{enumerate}
\item résoudre sur chacun des intervalles où $\alpha$ ne s'annule pas (deux constantes $K_-$ et $K_+$) ;
\item une solution sur $I$ tout entier est \emph{continue} en $t_0$ : égalité des limites à gauche et à droite ;
\item elle est \emph{dérivable} en $t_0$ : égalité des limites des taux d'accroissement (ou des dérivées, si elles ont une limite) ;
\item vérifier que l'équation est satisfaite en $t_0$ lui-même.
\end{enumerate}
Le nombre de solutions sur $I$ peut être $0$, $1$ ou infini : le théorème de Cauchy ne s'applique pas en $t_0$.
\end{methode}

\begin{exemple}[trois situations]
\begin{itemize}
\item $ty'+y=1$ : sur $\R_+^*$ et $\R_-^*$, $y=1+\frac{K_\pm}t$. La continuité en $0$ impose $K_+=K_-=0$ :
\emph{une seule} solution sur $\R$, $y=1$.
\item $ty'-y=t^2$ : sur chaque demi-droite, $y=t^2+K_\pm t$ ; ces fonctions tendent vers $0$ en $0$, et le taux d'accroissement
$\frac{y(t)}t=t+K_\pm$ tend vers $K_\pm$. La dérivabilité impose $K_+=K_-$ : \emph{une infinité} de solutions, $y=t^2+Kt$.
\item $ty'=1$ : $y=\ln|t|+K_\pm$ n'a pas de limite finie en $0$ : \emph{aucune} solution sur $\R$.
\end{itemize}
\end{exemple}

\fi

\begin{exo}{ed-o1}{1}{solutions particulières}
Résoudre sur $\R$ : (a) $y'+y=\sin t$ ; (b) $y'-2y=e^t\sin(2t)$ ; (c) $y'+y=te^{-t}$.
\end{exo}
\begin{exo}{ed-cauchy1}{2}{problème de Cauchy}
Résoudre sur $\left]-\frac\pi2,\frac\pi2\right[$ le problème $\cos(t)\,y'+\sin(t)\,y=1$, $y(0)=2$.
\end{exo}
\begin{exo}{ed-integ}{2}{équation intégrale}
Déterminer les fonctions $f$ continues sur $\R$ telles que, pour tout réel $x$, $f(x)-\displaystyle\int_0^x tf(t)\,dt=1+x$.
\end{exo}
\ifPT
\begin{exo}{ed-recol}{3}{recollement}
On considère $(E)$ : $|t|\,y'-y=t^2$.
\begin{enumerate}
\item Résoudre $(E)$ sur $\R_+^*$ et sur $\R_-^*$.
\item Déterminer les solutions de $(E)$ sur $\R$.
\end{enumerate}
\end{exo}
\fi

% ---------------------------------------------------------------------
\subsection{Second ordre à coefficients constants}\label{ss:constants}

On considère $(E)$ : $y''+ay'+by=c(t)$ avec $a,b\in\K$ constants et $c$ continue sur $I$. L'\emph{équation caractéristique}
est $r^2+ar+b=0$ ; elle vient de ce que $t\mapsto e^{rt}$ est solution de $(H)$ si et seulement si $r^2+ar+b=0$.

\begin{theoreme}[équation homogène]\label{t:caract}
\begin{itemize}
\item \emph{Cas complexe} ($\K=\C$). Si l'équation caractéristique a deux racines distinctes $r_1,r_2$, les solutions sont
$t\mapsto Ae^{r_1t}+Be^{r_2t}$ ; si elle a une racine double $r$, ce sont $t\mapsto(A+Bt)e^{rt}$ ($A,B\in\C$).
\item \emph{Cas réel} ($a,b\in\R$, solutions réelles). Discriminant $\Delta=a^2-4b$ :
\begin{itemize}
\item $\Delta>0$, racines réelles $r_1\neq r_2$ : $y=Ae^{r_1t}+Be^{r_2t}$ ;
\item $\Delta=0$, racine double $r$ : $y=(A+Bt)e^{rt}$ ;
\item $\Delta<0$, racines $\alpha\pm i\beta$ ($\beta>0$) : $y=e^{\alpha t}\big(A\cos(\beta t)+B\sin(\beta t)\big)$,
\end{itemize}
avec $A,B\in\R$.
\end{itemize}
\end{theoreme}

\begin{proof}
Soit $r$ une racine et $y=z\,e^{rt}$. Alors $y''+ay'+by=\big(z''+(2r+a)z'\big)e^{rt}$ (le terme en $z$ disparaît car
$r^2+ar+b=0$). Donc $y$ est solution si et seulement si $z'$ est solution de $w'+(2r+a)w=0$. Si $r$ est double, $2r+a=0$ :
$z''=0$, $z=A+Bt$. Sinon $2r+a=r-r_2\neq0$ : $z'=Ce^{(r_2-r)t}$, d'où $z=A+Be^{(r_2-r)t}$. Dans le cas réel avec $\Delta<0$,
on prend les parties réelle et imaginaire des solutions complexes $e^{(\alpha\pm i\beta)t}$, qui sont deux solutions réelles
non proportionnelles d'un plan vectoriel.
\end{proof}

\begin{methode}{Second membre $P(t)e^{mt}$}\label{m:part2}
Pour $y''+ay'+by=P(t)e^{mt}$ ($P$ polynôme, $m\in\K$), chercher $y=Q(t)e^{mt}$ avec $Q$ polynôme et
\[\deg Q=\begin{cases}\deg P&\text{si $m$ n'est pas racine de l'équation caractéristique},\\\deg P+1&\text{si $m$ est racine simple},\\\deg P+2&\text{si $m$ est racine double}.\end{cases}\]
En reportant, on obtient $Q''+(2m+a)Q'+(m^2+am+b)Q=P$ : on identifie les coefficients.
Cas particuliers : second membre polynomial ($m=0$) ; $\cos(\omega t)$, $\sin(\omega t)$, $e^{\alpha t}\cos(\omega t)$ : passer par
$e^{(\alpha+i\omega)t}$ et prendre la partie réelle ou imaginaire.
\end{methode}

\begin{astuce}
La formule $Q''+(2m+a)Q'+(m^2+am+b)Q=P$ évite de dériver deux fois un produit. Quand $m$ est racine double,
$m^2+am+b=0$ et $2m+a=0$ : il reste $Q''=P$, il suffit d'intégrer deux fois.
\end{astuce}

\begin{exemple}[superposition]
$y''-3y'+2y=2t+3+(6t+1)e^{-t}+e^t$. Racines $1$ et $2$ : $S_H=\{Ae^t+Be^{2t}\}$. On découpe le second membre :
\begin{itemize}
\item $2t+3$ ($m=0$, non racine) : $y_1=\alpha t+\beta$, $-3\alpha+2(\alpha t+\beta)=2t+3$, $y_1=t+3$ ;
\item $(6t+1)e^{-t}$ ($m=-1$, non racine) : $Q''-5Q'+6Q=6t+1$ avec $Q=\alpha t+\beta$ : $y_2=(t+1)e^{-t}$ ;
\item $e^t$ ($m=1$, racine simple) : $Q''-Q'=1$, $Q=-t$ : $y_3=-te^t$.
\end{itemize}
Solutions : $y=t+3+(t+1)e^{-t}-te^t+Ae^t+Be^{2t}$.
\end{exemple}

\begin{exemple}[résonance]
$y''+y=\sin t$. Ici $i$ est racine simple de $r^2+1$ : on cherche $z=Q(t)e^{it}$ pour $z''+z=e^{it}$, avec $\deg Q=1$.
$Q''+2iQ'=1$ donne $Q=-\frac{i}2t$, $z=-\frac{it}2e^{it}$, et $y_p=\operatorname{Im}z=-\frac t2\cos t$.
Solutions : $y=-\frac t2\cos t+A\cos t+B\sin t$. L'amplitude de $y_p$ croît sans limite : c'est la résonance de l'oscillateur
non amorti excité à sa propre pulsation.
\end{exemple}

\begin{enpratique}
L'oscillateur amorti $y''+2\lambda y'+\omega^2y=0$ ($\lambda,\omega>0$) a pour discriminant réduit $\lambda^2-\omega^2$ :
régime \emph{apériodique} si $\lambda>\omega$ (deux exponentielles décroissantes), \emph{critique} si $\lambda=\omega$,
\emph{pseudo-périodique} si $\lambda<\omega$ (oscillations de pulsation $\sqrt{\omega^2-\lambda^2}$ amorties par $e^{-\lambda t}$).
\end{enpratique}

\ifPT
\begin{methode}{Variation de la constante au second ordre}\label{m:var2}
Si le second membre n'a pas la forme $P(t)e^{mt}$, prendre une solution $\varphi$ de $(H)$ qui ne s'annule pas sur $I$
(par exemple $e^{rt}$) et chercher $y=\lambda(t)\varphi(t)$. En reportant, le terme en $\lambda$ disparaît et on obtient une
équation du \emph{premier ordre} en $\lambda'$ : la résoudre, puis intégrer.
\end{methode}

\begin{exemple}
$y''-2y'+y=\frac{2te^t}{1+t^2}$. Racine double $1$, $\varphi(t)=e^t$. Avec $y=\lambda e^t$ : $y''-2y'+y=\lambda''e^t$, donc
$\lambda''=\frac{2t}{1+t^2}$, $\lambda'=\ln(1+t^2)+C$, et par intégration par parties
$\lambda=t\ln(1+t^2)-2t+2\arctan t+Ct+D$. Solutions : $y=\big(t\ln(1+t^2)-2t+2\arctan t\big)e^t+(Ct+D)e^t$.
\end{exemple}
\else
\begin{attention}
Pour un second membre qui n'est pas de l'une de ces formes, il n'y a pas de méthode générale en première année : l'énoncé
guide le calcul (voir par exemple le problème du TD sur les solutions périodiques d'un oscillateur forcé).
\end{attention}
\fi

\begin{exo}{ed-c1}{1}{équations homogènes}
Résoudre sur $\R$ : (a) $y''-3y'+2y=0$ ; (b) $y''+2y'+y=0$ ; (c) $y''+y'+y=0$ (solutions réelles).
\end{exo}
\begin{exo}{ed-c2}{2}{seconds membres}
Résoudre sur $\R$ : (a) $y''+4y'+4y=\sin t$ ; (b) $y''-2y'+y=e^{mt}$, selon le réel $m$.
\end{exo}
\begin{exo}{ed-c3}{2}{résonance et problème de Cauchy}
Déterminer la solution de $y''+y=\sin t$ telle que $y(0)=0$ et $y'(0)=0$. Est-elle bornée sur $\R$ ?
\end{exo}

\ifPT
% ---------------------------------------------------------------------
\subsection{Second ordre à coefficients continus}\label{ss:continus}

On considère $(E)$ : $\alpha(t)y''+\beta(t)y'+\gamma(t)y=\delta(t)$, avec des fonctions continues et $\alpha$ \emph{qui ne
s'annule pas} sur $I$. Il n'y a plus de formule générale ; le théorème de Cauchy et la structure restent vrais
(théorème~\ref{t:cauchy}, corollaire~\ref{c:dim}) : il suffit de trouver \emph{deux solutions non proportionnelles} de $(H)$, puis
une solution particulière. L'énoncé indique en général la méthode.

\begin{methode}{Chercher une solution polynomiale ou exponentielle}\label{m:poly}
\emph{Polynomiale :} on pose $y=\sum_{k=0}^na_kt^k$ avec $a_n\neq0$ et l'on écrit que le coefficient de plus haut degré est nul :
cela fixe $n$. On identifie ensuite les autres coefficients.\\
\emph{Exponentielle :} on reporte $y=e^{rt}$ ; l'égalité doit avoir lieu pour tout $t$, d'où des conditions sur $r$.
\end{methode}

\begin{exemple}
$(2t+1)y''+(4t-2)y'-8y=0$. Pour $y$ de degré $n$, le coefficient de $t^n$ est $(4n-8)a_n$ : $n=2$. Avec $y=at^2+bt+c$ on trouve
$b=0$, $c=\frac a4$ : $y_1=4t^2+1$. Avec $y=e^{rt}$ : $(2t+1)r^2+(4t-2)r-8=0$ pour tout $t$, soit $2r^2+4r=0$ et $r^2-2r-8=0$ :
$r=-2$, $y_2=e^{-2t}$. Sur tout intervalle ne contenant pas $-\frac12$, $y_1$ et $y_2$ ne sont pas proportionnelles :
$S_H=\{A(4t^2+1)+Be^{-2t}\}$.
\end{exemple}

\begin{methode}{Une solution $\varphi$ de $(H)$ ne s'annulant pas est connue}\label{m:abaisse}
Chercher $y=\lambda(t)\varphi(t)$ (programme). En reportant, le terme en $\lambda$ disparaît (car $\varphi$ est solution) :
\[\alpha\varphi\,\lambda''+(2\alpha\varphi'+\beta\varphi)\,\lambda'=\delta,\]
équation du \emph{premier ordre} en $\mu=\lambda'$. On la résout (§\ref{ss:ordre1}), puis on intègre $\mu$. Les deux constantes
d'intégration donnent la solution générale.
\end{methode}

\begin{exemple}[d'après un TD]
$t^2y''-2y=3t^2$ sur $\R_+^*$. Un polynôme solution de $(H)$ a un coefficient dominant vérifiant $n(n-1)-2=0$ : $n=2$, et
$\varphi(t)=t^2$ convient. Avec $y=\lambda t^2$ : $t^4\lambda''+4t^3\lambda'=3t^2$, soit $\mu'+\frac4t\mu=\frac3{t^2}$.
L'homogène donne $\frac K{t^4}$, la variation de la constante $\mu=\frac1t$ : $\mu=\frac1t+\frac K{t^4}$, puis
$\lambda=\ln t-\frac K{3t^3}+D$. Solutions : $y=t^2\ln t+\frac at+bt^2$, $(a,b)\in\R^2$.
\end{exemple}

\begin{methode}{Changement de variable}\label{m:chgt}
L'énoncé propose $t=\varphi(x)$ avec $\varphi$ bijective et $\mathcal C^2$, de réciproque $\mathcal C^2$. On pose
$z(x)=y(\varphi(x))$, on calcule $z'(x)=\varphi'(x)y'(\varphi(x))$ et $z''(x)$, et l'on traduit l'équation en $z$, qui doit
devenir une équation connue (souvent à coefficients constants). On revient enfin à $y=z\circ\varphi^{-1}$.
\end{methode}

\begin{exemple}
$(1-t^2)y''-ty'+y=0$ sur $]-1,1[$, avec $t=\cos x$, $x\in\,]0,\pi[$. $z'(x)=-\sin x\,y'(\cos x)$ et
$z''(x)=\sin^2x\,y''(\cos x)-\cos x\,y'(\cos x)$. L'équation devient $z''+z=0$ : $z=A\cos x+B\sin x$, et comme
$\sin(\arccos t)=\sqrt{1-t^2}$, $y=At+B\sqrt{1-t^2}$.
\end{exemple}

\begin{plusloin}
L'équation d'Euler $t^2y''+aty'+by=0$ sur $\R_+^*$ se ramène à coefficients constants par $t=e^x$ :
$z''+(a-1)z'+bz=0$. Par exemple $t^2y''-ty'+y=0$ donne $z''-2z'+z=0$, $z=(A+Bx)e^x$, d'où $y=t(A+B\ln t)$.
\end{plusloin}

\begin{methode}{Solution développable en série entière}\label{m:se}
\emph{Analyse :} supposer $y(t)=\sum a_nt^n$ de rayon $R>0$, dériver terme à terme, regrouper par puissance de $t$ (décaler
les indices), et identifier (unicité du développement) : on obtient une relation de récurrence sur les $a_n$.
\emph{Synthèse :} calculer les $a_n$, déterminer le rayon de la série obtenue, et conclure que sa somme est solution sur
$]-R,R[$ ; reconnaître si possible une fonction usuelle.
\end{methode}

\begin{exemple}
$(t^2+t)y''+(3t+1)y'+y=0$. Le coefficient de $t^n$ est $n(n-1)a_n+(n+1)na_{n+1}+3na_n+(n+1)a_{n+1}+a_n=(n+1)^2(a_n+a_{n+1})$.
Donc $a_{n+1}=-a_n$, $a_n=(-1)^na_0$, de rayon $1$ : $y=\frac{a_0}{1+t}$ sur $]-1,1[$.
\end{exemple}

\begin{attention}
Ces méthodes donnent les solutions sur un intervalle où le coefficient de $y''$ ne s'annule pas. Une série entière n'est
solution que sur son intervalle de convergence ; elle ne donne en général qu'une partie de $S_H$.
\end{attention}

\begin{methode}{Développer une fonction en série entière grâce à une équation différentielle}\label{m:dse}
Le programme demande aussi le chemin inverse : pour développer une fonction $f$ connue,
\begin{enumerate}
\item trouver une équation différentielle linéaire, avec conditions initiales, dont $f$ est solution sur un intervalle $I$ contenant $0$ ;
\item chercher la solution développable en série entière de ce problème (méthode précédente) : récurrence sur les $a_n$, rayon $R>0$ ;
\item conclure par l'unicité du théorème de Cauchy : $f$ et la somme de la série sont solutions du même problème sur $]-R,R[\cap I$,
donc elles sont égales.
\end{enumerate}
\end{methode}

\begin{exemple}[la série du binôme]
Soit $\alpha\in\R\setminus\N$ et $f(x)=(1+x)^\alpha$ sur $]-1,1[$ : $(1+x)f'(x)=\alpha f(x)$ et $f(0)=1$. Si $y=\sum a_nx^n$ est solution
avec $a_0=1$, le coefficient de $x^n$ donne $(n+1)a_{n+1}+na_n=\alpha a_n$, soit $a_{n+1}=\frac{\alpha-n}{n+1}a_n$, d'où
\[a_n=\frac{\alpha(\alpha-1)\cdots(\alpha-n+1)}{n!}.\]
Comme $\left|\frac{a_{n+1}}{a_n}\right|=\frac{|\alpha-n|}{n+1}\to1$, le rayon vaut $1$. La somme et $f$ sont solutions sur $]-1,1[$ du même
problème de Cauchy (premier ordre, $1+x\neq0$) : elles sont égales. Ainsi $(1+x)^\alpha=\sum_{n\geq0}\frac{\alpha(\alpha-1)\cdots(\alpha-n+1)}{n!}x^n$ pour $|x|<1$.
\end{exemple}

\ifPT
\begin{exo}{ed-poly}{2}{solution polynomiale}
On considère $(E)$ : $(1+t^2)y''+ty'-y=0$ sur $\R$.
\begin{enumerate}
\item Déterminer les solutions polynomiales de $(E)$.
\item Résoudre $(E)$ en posant $t=\sh x$.
\item Déterminer la solution telle que $y(0)=1$ et $y'(0)=1$.
\end{enumerate}
\end{exo}
\fi
\ifPT
\begin{exo}{ed-abaisse}{3}{une solution connue}
On considère $(E)$ : $(t+1)y''-(t+2)y'+y=0$ sur $]-1,+\infty[$.
\begin{enumerate}
\item Déterminer une solution de la forme $t\mapsto e^{rt}$.
\item En déduire toutes les solutions de $(E)$.
\end{enumerate}
\end{exo}
\fi
\ifPT
\begin{exo}{ed-se}{3}{série entière}
Déterminer la solution développable en série entière du problème $y''-2ty'-2y=0$, $y(0)=1$, $y'(0)=0$, et l'exprimer à l'aide
des fonctions usuelles.
\end{exo}
\fi

\fi
% ---------------------------------------------------------------------
\subsection{Quelle méthode ? Synthèse}\label{ss:synthese}

\ifPT
\begin{center}
\renewcommand{\arraystretch}{1.35}
\small
\begin{tabular}{p{0.38\linewidth}p{0.55\linewidth}}
\hline
\textbf{Équation} & \textbf{Méthode}\\\hline
$y'+a(t)y=b(t)$ & homogène $Ke^{-A}$ ; particulière évidente, de la forme du second membre, ou variation de la constante (méthode~\ref{m:part1})\\
$\alpha(t)y'+\dots$ avec $\alpha$ qui s'annule & résoudre sur chaque intervalle, puis recoller (méthode~\ref{m:recollement})\\
$y''+ay'+by=c(t)$, $a,b$ constants & équation caractéristique ; second membre $P(t)e^{mt}$ (méthode~\ref{m:part2}) ; sinon variation de la constante (méthode~\ref{m:var2})\\
$\alpha(t)y''+\beta(t)y'+\gamma(t)y=\delta(t)$ & suivre l'indication : solution polynomiale ou exponentielle (\ref{m:poly}), solution connue (\ref{m:abaisse}), changement de variable (\ref{m:chgt}), série entière (\ref{m:se})\\
inconnue sous une intégrale & dériver (méthode~\ref{m:integrale})\\\hline
\end{tabular}
\end{center}
\else
\begin{center}
\renewcommand{\arraystretch}{1.35}
\small
\begin{tabular}{p{0.38\linewidth}p{0.55\linewidth}}
\hline
\textbf{Équation} & \textbf{Méthode}\\\hline
$y'+a(t)y=b(t)$ & homogène $Ke^{-A}$ ; particulière évidente, de la forme du second membre, ou variation de la constante (méthode~\ref{m:part1})\\
$y''+ay'+by=c(t)$, $a,b$ constants & équation caractéristique ; second membre $P(t)e^{mt}$ (méthode~\ref{m:part2}) ; superposition si le second membre est une somme\\
$\alpha(t)y'+\beta(t)y=\gamma(t)$ & normaliser sur chaque intervalle où $\alpha$ ne s'annule pas\\
inconnue sous une intégrale & dériver (méthode~\ref{m:integrale})\\\hline
\end{tabular}
\end{center}
\fi

\begin{attention}\textbf{Pièges fréquents.}
\begin{itemize}[leftmargin=1.1em]
\item Oublier de normaliser avant de prendre la primitive, ou se tromper de signe dans $e^{-A}$.
\item Oublier la résonance : si $m$ est racine, la forme $Ce^{mt}$ ne peut pas marcher.
\item Prendre la partie réelle au lieu de la partie imaginaire (second membre en $\sin$).
\item Fixer les constantes avec les conditions initiales \emph{avant} d'avoir ajouté la solution particulière.
\item Diviser par un coefficient qui s'annule sans le signaler\ifPT, ou affirmer l'unicité d'une solution recollée sans vérifier\fi.
\end{itemize}
\end{attention}

\begin{essentiel}
\begin{itemize}[leftmargin=1.1em,itemsep=2pt]
\item $S_E=y_p+S_H$ ; $S_H$ est une droite (ordre 1) ou un plan (ordre 2)\ifPT{} dès que le coefficient dominant ne s'annule pas\fi.
\item Cauchy : une condition $y(t_0)$ à l'ordre 1, deux conditions $y(t_0),y'(t_0)$ à l'ordre 2, déterminent une unique solution.
\item Ordre 1 : $y=Ke^{-A(t)}$, puis variation de la constante $\lambda'=b\,e^{A}$.
\item Coefficients constants : $r^2+ar+b=0$ ; second membre $P(t)e^{mt}$, degré augmenté de la multiplicité de $m$.
\ifPT
\item Coefficients continus : deux solutions non proportionnelles de $(H)$ suffisent ; une solution connue $\varphi$ ramène au
premier ordre par $y=\lambda\varphi$.
\fi
\end{itemize}
\end{essentiel}

% =====================================================================
%                              CORRIGÉS
% =====================================================================
\clearpage
\renewcommand{\ongletcouleur}{Ocre}
\section{Corrigés}\label{sec:corriges}

\begin{corrige}{ed-o1}
\q{(a)} Homogène : $Ke^{-t}$. Par $z'+z=e^{it}$, $z=\frac{e^{it}}{1+i}=\frac{(1-i)e^{it}}2$, d'où $y_p=\operatorname{Im}z=\frac{\sin t-\cos t}2$.
Solutions : $y=\frac{\sin t-\cos t}2+Ke^{-t}$.

\q{(b)} Homogène : $Ke^{2t}$. Le second membre est $\operatorname{Im}\big(e^{(1+2i)t}\big)$ et $1+2i\neq2$ : on cherche $z=\nu e^{(1+2i)t}$,
$\nu(1+2i-2)=1$, $\nu=\frac1{-1+2i}=\frac{-1-2i}5$. Donc $y_p=\operatorname{Im}z=-\frac{e^t}5\big(\sin2t+2\cos2t\big)$ et
$y=Ke^{2t}-\frac{e^t}5(\sin2t+2\cos2t)$.

\q{(c)} $m=-1=-a$ : résonance, on cherche $y=Q(t)e^{-t}$ avec $\deg Q=2$. En reportant : $Q'=t$, $Q=\frac{t^2}2$.
Solutions : $y=\left(\frac{t^2}2+K\right)e^{-t}$.
\piege{en (c), chercher $(\alpha t+\beta)e^{-t}$ échoue : le terme $\beta e^{-t}$ est solution de l'homogène.}
\end{corrige}

\begin{corrige}{ed-cauchy1}
Sur $\left]-\frac\pi2,\frac\pi2\right[$, $\cos t\neq0$ : $y'+\tan(t)\,y=\frac1{\cos t}$. Une primitive de $\tan$ est $-\ln(\cos t)$, d'où
l'homogène $y=Ke^{\ln\cos t}=K\cos t$. Variation de la constante : $y=\lambda\cos t$ donne $\lambda'\cos t=\frac1{\cos t}$,
$\lambda=\tan t+K$. Solutions : $y=\sin t+K\cos t$ ; $y(0)=K=2$. La solution est $y=\sin t+2\cos t$.
\gain{une solution particulière évidente ($\sin t$) se vérifie en une ligne : $\cos t\cos t+\sin t\sin t=1$.}
\end{corrige}

\begin{corrige}{ed-integ}
\choix{l'inconnue est sous une intégrale de borne variable : on dérive (méthode~\ref{m:integrale}).}
Si $f$ convient, $f(x)=1+x+\int_0^xtf(t)\,dt$ est dérivable, $f'(x)=1+xf(x)$, et $f(0)=1$. Réciproquement, une solution de ce
problème de Cauchy, intégrée de $0$ à $x$, vérifie l'équation de départ. On résout $y'-xy=1$ : l'homogène donne $Ke^{x^2/2}$ ;
la variation de la constante donne $\lambda'=e^{-x^2/2}$, sans primitive usuelle. On écrit donc la solution du problème de Cauchy
sous forme intégrale :
\[f(x)=e^{x^2/2}\left(1+\int_0^xe^{-u^2/2}\,du\right).\]
\retenir{une primitive sans expression usuelle s'écrit $\int_{x_0}^x$ ; la borne $x_0$ est choisie pour que la condition initiale soit simple.}
\end{corrige}

\ifPT
\begin{corrige}{ed-recol}
\q{1.} Sur $\R_+^*$ : $ty'-y=t^2$, soit $y'-\frac yt=t$. Homogène : $Kt$ ; $y=t^2$ convient. Solutions : $y=t^2+K_+t$.
Sur $\R_-^*$ : $-ty'-y=t^2$, soit $y'+\frac yt=-t$. Homogène : $\frac Kt$ ; on cherche $y=\alpha t^2$ : $2\alpha t+\alpha t=-t$,
$\alpha=-\frac13$. Solutions : $y=-\frac{t^2}3+\frac{K_-}t$.

\q{2.} Soit $y$ solution sur $\R$. Elle est continue en $0$ : la limite à gauche n'est finie que si $K_-=0$, et alors elle vaut
$0$ ; donc $y(0)=0$. À droite, $\frac{y(t)}t=t+K_+\to K_+$ ; à gauche, $\frac{y(t)}t=-\frac t3\to0$. La dérivabilité en $0$
impose $K_+=0$. Réciproquement, la fonction $y$ égale à $t^2$ sur $\R_+$ et à $-\frac{t^2}3$ sur $\R_-$ est dérivable en $0$
($y'(0)=0$), et l'équation est vérifiée en $0$ ($0-0=0$). Il y a donc une unique solution sur $\R$.
\retenir{continuité, puis dérivabilité, puis vérification de l'équation au point de raccord.}
\end{corrige}
\fi

\begin{corrige}{ed-c1}
(a) $r^2-3r+2=(r-1)(r-2)$ : $y=Ae^t+Be^{2t}$. (b) $(r+1)^2$ : $y=(A+Bt)e^{-t}$.
(c) $r=-\frac12\pm i\frac{\sqrt3}2$ : $y=e^{-t/2}\left(A\cos\frac{\sqrt3}2t+B\sin\frac{\sqrt3}2t\right)$.
\end{corrige}

\begin{corrige}{ed-c2}
\q{(a)} Racine double $-2$ : $S_H=\{(A+Bt)e^{-2t}\}$. Pour $z''+4z'+4z=e^{it}$, $i$ n'est pas racine : $z=\frac{e^{it}}{(i+2)^2}=\frac{e^{it}}{3+4i}=\frac{(3-4i)e^{it}}{25}$,
d'où $y_p=\operatorname{Im}z=\frac{3\sin t-4\cos t}{25}$ et $y=(A+Bt)e^{-2t}+\frac{3\sin t-4\cos t}{25}$.

\q{(b)} Racine double $1$ : $S_H=\{(A+Bt)e^t\}$. Si $m\neq1$, $m$ n'est pas racine : $y_p=\frac{e^{mt}}{(m-1)^2}$. Si $m=1$, racine
double : $Q''=1$, $y_p=\frac{t^2}2e^t$.
\retenir{l'expression de $y_p$ change de forme exactement quand $m$ devient racine : discuter selon la multiplicité.}
\end{corrige}

\begin{corrige}{ed-c3}
D'après l'exemple de résonance, $y=-\frac t2\cos t+A\cos t+B\sin t$. $y(0)=A=0$ ; $y'(t)=-\frac12\cos t+\frac t2\sin t+B\cos t$,
$y'(0)=B-\frac12=0$. La solution est $y=\frac{\sin t-t\cos t}2$. Elle n'est pas bornée : $y(2k\pi)=-k\pi\to-\infty$.
\end{corrige}

\ifPT
\begin{corrige}{ed-poly}
\q{1.} Pour $y$ de degré $n$, le coefficient de $t^n$ est $\big(n(n-1)+n-1\big)a_n=(n^2-1)a_n$ : $n=1$ (le cas constant donne
$y=0$). Avec $y=at+b$ : $at-at-b=0$, donc $b=0$. Solutions polynomiales : $y=at$.

\q{2.} $t=\sh x$ est une bijection $\mathcal C^2$ de $\R$ sur $\R$, de réciproque $\mathcal C^2$. Avec $z(x)=y(\sh x)$ :
$z'=\ch x\,y'(\sh x)$, $z''=\ch^2x\,y''(\sh x)+\sh x\,y'(\sh x)$. Comme $1+\sh^2x=\ch^2x$, l'équation devient $z''-z=0$ :
$z=A\sh x+B\ch x$, d'où $y=At+B\sqrt{1+t^2}$.

\q{3.} $y(0)=B=1$ et $y'(0)=A=1$ : $y=t+\sqrt{1+t^2}$.
\retenir{le changement de variable proposé transforme $1+t^2$ en $\ch^2x$ : l'équation devient à coefficients constants.}
\end{corrige}
\fi

\ifPT
\begin{corrige}{ed-abaisse}
\q{1.} $y=e^{rt}$ : $(t+1)r^2-(t+2)r+1=0$ pour tout $t$, soit $r^2-r=0$ et $r^2-2r+1=0$ : $r=1$. Donc $\varphi(t)=e^t$.

\q{2.} $\varphi$ ne s'annule pas : on pose $y=\lambda e^t$. Alors $y'=(\lambda'+\lambda)e^t$, $y''=(\lambda''+2\lambda'+\lambda)e^t$ et
l'équation devient $(t+1)\lambda''+t\lambda'=0$, soit $\mu'=-\frac t{t+1}\mu=\left(-1+\frac1{t+1}\right)\mu$ pour $\mu=\lambda'$.
D'où $\mu=K(t+1)e^{-t}$, puis $\lambda=-K(t+2)e^{-t}+D$ (vérifier : $\big((t+2)e^{-t}\big)'=-(t+1)e^{-t}$).
Solutions : $y=a(t+2)+be^t$, $(a,b)\in\R^2$.
\gain{on peut vérifier directement que $t+2$ est solution : $0-(t+2)+(t+2)=0$.}
\end{corrige}
\fi

\ifPT
\begin{corrige}{ed-se}
\choix{analyse-synthèse : relation de récurrence par identification, puis rayon de convergence.}
Avec $y=\sum a_nt^n$ : $\sum(n+2)(n+1)a_{n+2}t^n-\sum2na_nt^n-\sum2a_nt^n=0$, d'où $(n+2)(n+1)a_{n+2}=2(n+1)a_n$, soit
$a_{n+2}=\frac2{n+2}a_n$. Avec $a_1=y'(0)=0$, les termes impairs sont nuls ; $a_0=1$ et $a_{2p}=\frac2{2p}a_{2p-2}=\frac1{p!}$.
Donc $y=\sum\frac{t^{2p}}{p!}=e^{t^2}$, de rayon infini. Vérification : $y'=2te^{t^2}$, $y''=(2+4t^2)e^{t^2}$ et
$y''-2ty'-2y=0$.
\end{corrige}
\fi

\end{document}
